\begin{corollary}
Let $Q$ be a quasi-order.  Then $Q$ is bqo if and only if for every
locally constant $\varphi:\mathsf{IncSeq}\to Q$ and every
$g\in\mathsf{IncSeq}$ there exists $f\in\mathsf{IncSeq}$ such that
\[
\varphi(f)\leq\varphi(f\circ g).
\]
\end{corollary}
\begin{proof}
Assume first that $Q$ is bqo, and let $\varphi$ be locally constant and
$g\in\mathsf{IncSeq}$.  If $g=\mathsf{IncSeqId}$, choose any $f$, for
example $f=\mathsf{IncSeqId}$.  By the identity definition in
\noderef{IncSeqId}, and since $Q$ is a quasi-order, reflexivity
gives $\varphi(f)\leq\varphi(f\circ g)$, because $f\circ
\mathsf{IncSeqId}=f$.

Suppose instead that $g\neq\mathsf{IncSeqId}$.  The quasi-order
specialization of \noderef{GBetterRelIff} says that $Q$ is
$g$-better.  Use the first clause of \noderef{GBetterRel} with
$\varphi$.  It supplies $f$ and a continuous homomorphism $H$ for
$(\mathsf{IncSeq},\mathsf{RightComp}(g))$ such that
$H(x)=\varphi(f\circ x)$.  Apply the homomorphism relation from
\noderef{ContinuousRelHom} at $x=\mathsf{IncSeqId}$.  Substituting the displayed
formula for $H$ and using the definitions of \noderef{RightComp} and
composition gives
\[
 r(\varphi(f),\varphi(f\circ g)),
\]
as required.

Conversely, assume the displayed property for every locally constant
$\varphi$ and every $g$.  Let $h$ be locally constant.  If $h$ were bad,
then \noderef{BadMultiSequence} would give
\[
\neg r(h(f),h(\mathsf{ShiftMap}(f)))
\]
for every $f$.  Apply the assumed property to $\varphi=h$ and
$g=\mathsf{SuccSeq}$.  It yields an $f$ with
$r(h(f),h(\mathsf{RightComp}(\mathsf{SuccSeq})(f)))$.  By
\noderef{ShiftMap} and \noderef{SuccSeq}, the second argument is
$h(\mathsf{ShiftMap}(f))$,
contradicting badness.  Hence no locally constant map is bad, and
\noderef{Bqo} proves that $Q$ is bqo.
\end{proof}
